% ==============================================================================
% SECCIÓN 5: JUSTIFICACIÓN DE LA PROPUESTA
% ==============================================================================
\newpage
\section{Justificación de la Propuesta}

Las decisiones siguientes son la justificación inicial del equipo para acotar la primera versión. Se apoyan en el propósito del proyecto —apoyar la organización y la autonomía del alumnado PTVAL— y deben contrastarse con las solicitudes originales del centro y con la información recogida en reuniones. No se presentan como necesidades confirmadas hasta completar esa trazabilidad.

\subsection{Soluciones relacionadas y comparativa}
Como revisión preliminar se comparan dos herramientas digitales y un recurso de apoyos visuales descritos por sus entidades responsables. La comparación identifica solapamientos y el foco propuesto para AUTOVIDA; no afirma que las soluciones existentes carezcan de funciones que no aparezcan en sus descripciones públicas.

\begin{center}
\small
\begin{tabularx}{\textwidth}{|>{\raggedright\arraybackslash}p{3.1cm}|>{\raggedright\arraybackslash}X|>{\raggedright\arraybackslash}p{5.1cm}|}
\hline
\rowcolor{tableHeader}
Referencia & Capacidades descritas públicamente & Relación con AUTOVIDA \\
\hline
\href{https://fundacionorange.es/aplicaciones/in-tic-agenda/}{In-TIC Agenda (Fundación Orange / IMEDIR)} & Permite crear agendas diarias, semanales y libres con actividades asociadas a secuencias de pictogramas, configurables por profesionales o familiares. & Coincide con la organización visual y secuenciación de actividades. AUTOVIDA propone como alcance propio la asignación de tareas por perfil, el registro de pasos completados y un temporizador vinculado a una actividad. \\
\hline
\href{https://fundacionorange.es/aplicaciones/dia-a-dia/}{Día a Día (Fundación Orange)} & Diario visual para guardar y revisar actividades mediante imágenes, fotos, vídeos y textos; permite situarlas en momentos del día y asociarlas con personas y lugares. & Su descripción se centra en el diario visual de actividades. AUTOVIDA plantea centrarse en la agenda de actividades pendientes asignadas por el personal y en acompañar su realización. \\
\hline
\href{https://aulaabierta.arasaac.org/materiales-caa-agendas-visuales}{Agendas visuales (ARASAAC)} & Reúne materiales de agendas visuales con pictogramas y explica su uso para evocar actividades previstas o realizadas, apoyar la anticipación y facilitar la comunicación entre centro y familia. & Aporta una referencia para los apoyos visuales. AUTOVIDA propone integrarlos en una aplicación con perfiles, tareas y seguimiento del avance; la selección de pictogramas y sus condiciones de uso se verificarán antes de incorporarlos. \\
\hline
\end{tabularx}
\end{center}

La revisión se basa en las descripciones oficiales consultadas el 8 de octubre de 2026. Antes de cerrar la propuesta se comprobarán las versiones disponibles, dispositivos compatibles, condiciones de uso y adecuación con el centro; esta tabla no sustituye a una prueba práctica de las herramientas.

\begin{center}
\small
\begin{tabularx}{\textwidth}{|>{\raggedright\arraybackslash}p{2.8cm}|>{\raggedright\arraybackslash}p{2.7cm}|>{\raggedright\arraybackslash}X|}
\hline
\rowcolor{tableHeader}
Decisión propuesta & Requisitos relacionados & Justificación y efecto en el alcance \\
\hline
Incluir agenda y actividades asignadas en el MVP. & RF-07, RF-10 y RF-21. & Responde directamente al propósito de organizar tareas cotidianas. Limitar la agenda a actividades asignadas a cada perfil permite probar el flujo principal sin introducir calendarios compartidos o integraciones externas. \\
\hline
Incluir instrucciones ordenadas y registro de avance. & RF-07, RF-10, RF-25, RF-27 y RF-28. & Dividir una actividad en pasos observables permite comprobar qué instrucciones se han presentado y completado. La redacción, extensión y apoyo visual de cada paso se validarán con el centro y las personas usuarias. \\
\hline
Incluir rutinas configurables por el personal. & RF-12 y RF-13. & Permite preparar secuencias recurrentes desde el centro y evita que cada alumno tenga que configurar por sí mismo el contenido de sus apoyos. La frecuencia y forma de repetición quedan por concretar. \\
\hline
Incluir un temporizador visual sencillo. & RF-15, RF-33 y RF-35. & Se propone como apoyo a actividades con duración definida. Se limita a iniciar una duración configurada y señalar visualmente su finalización; pausas, avisos sonoros y otras opciones quedan fuera hasta validar su necesidad. \\
\hline
Posponer módulos de reprografía, menús/dietas, visitas y mensajería. & Requisitos futuros por identificar. & Son flujos distintos del núcleo de agenda y tareas. Los menús/dietas y la mensajería pueden requerir decisiones adicionales sobre datos y permisos; reprografía y visitas dependen de procesos específicos del centro. Posponerlos reduce el alcance inicial sin cerrar su incorporación futura. \\
\hline
\end{tabularx}
\end{center}

\subsection{Criterios para confirmar o cambiar las decisiones}
El equipo revisará estas decisiones con el centro. Para cada una, anotará qué se ha solicitado, qué se acuerda (mantener, modificar o posponer) y el motivo. Si cambia un requisito, también se actualizarán su prioridad y esta justificación.

\begin{infobox}[title={Formato de Solicitud}]
  Decisión: [Acciones Tomadas]. Motivo: [Razon para el cambio]. Persona que valida y fecha: [Integrante Validador y Fecha].
\end{infobox}